\documentclass[12pt, a4paper]{article}
% ── Packages ──────────────────────────────────────────────────────────────────
\usepackage[a4paper, margin=2.5cm]{geometry}
\usepackage{graphicx}
\usepackage{hyperref}
\usepackage{setspace}
\usepackage{titlesec}
\usepackage{parskip}
\usepackage{enumitem}
\usepackage{booktabs}
\usepackage{amsmath}
\usepackage{natbib}
\usepackage{xcolor}
\usepackage{fancyhdr}
\usepackage{microtype}
\usepackage[utf8]{inputenc}
\usepackage[T1]{fontenc}
% ── Page style ────────────────────────────────────────────────────────────────
\pagestyle{fancy}
\fancyhf{}
\rhead{\small Exploring Art in Medicine through Knowledge Graphs}
\lhead{\small Micky Kotterer}
\cfoot{\thepage}
% ── Hyperlink colours ─────────────────────────────────────────────────────────
\hypersetup{
    colorlinks=true,
    linkcolor=black,
    citecolor=black,
    urlcolor=blue
}
% ── Line spacing ──────────────────────────────────────────────────────────────
\onehalfspacing
% ── Title ─────────────────────────────────────────────────────────────────────
\title{%
    \Large\textbf{Exploring Art in Medicine through Knowledge Graphs}\\[0.5em]
    \large Thesis Draft --- University of Amsterdam
}
\author{Micky Kotterer (13464167)\\[0.3em]
    \small Supervisor: Natallia Kokash}
\date{June 2026}
% ══════════════════════════════════════════════════════════════════════════════
\begin{document}
\maketitle
\thispagestyle{empty}
% ── Abstract ──────────────────────────────────────────────────────────────────
\begin{abstract}
\noindent
This thesis presents a graph-based exploration system for the Bagatelle collection, a curated set of
roughly 600 artworks related to art in medicine. Standard keyword search is a poor fit for this kind
of collection: its meaning is largely visual, symbolic, and historically layered in ways that plain
metadata cannot capture. The system built in this project extends the existing Bagatelle server,
which already provides multimodal embeddings and LLM-generated artwork descriptions, with a
knowledge graph layer that lets users navigate the collection by following semantic connections
between artworks rather than searching for terms.

Five graph generation strategies have been implemented and compared: text-based similarity using
LLM-generated descriptions, image-based similarity using CLIP visual embeddings, a combined
approach that fuses both, a query-targeted mode that keeps expansion anchored to the user's
original search intent, and an LLM-confirmed mode where a language model validates and explains
each proposed connection. An interactive interface allows users to build and expand the graph step
by step, with every edge carrying both a similarity score and a natural-language explanation of why
two artworks are connected.

Evaluation is planned through a side-by-side comparison study in which users and domain experts
rate which strategy produces more relevant, interesting, and navigable connections. An MD evaluator
with connections to the medical humanities community has been engaged for this phase. The thesis
contributes a working exploration prototype, a configurable multi-strategy graph architecture, and
a design for interpretable edges in a small but semantically rich cultural heritage collection.
\end{abstract}
\newpage
% ── Table of Contents ─────────────────────────────────────────────────────────
\tableofcontents
\newpage
% ══════════════════════════════════════════════════════════════════════════════
\section{Introduction}

\subsection{Digital Exploration of Cultural Collections}

Cultural heritage collections are increasingly accessible online; however, accessibility alone does
not equal meaningful exploration. Large digitised repositories have been built by museums, archives,
and digital libraries, but the dominant interaction model, a search box and a list of results, still
treats the collection as a flat inventory rather than a web of connected knowledge. Several
researchers have argued that effective access to cultural heritage requires moving beyond simple
retrieval and toward systems that support structured, multidimensional exploration of collections
\citep{shiri2015, ardissono2016}.

The Bagatelle collection examined in this project is a specific case of this broader challenge. It
contains roughly 600 artworks relating to art in medicine: works depicting medical conditions,
clinical settings, anatomical studies, and historical themes of care and disease. This is not a
collection that responds well to keyword search. A query for ``plague'' will not return artworks
that depict the social isolation of the sick, or the symbolic use of the crow, without explicit
metadata pointing to those concepts. The collection's richness lives in the relationships between
works, not in the labels assigned to them.

This gap between what a collection contains and what a user can discover through it is the
motivating problem of this thesis.

\subsection{Limitations of Keyword-Based and Metadata-Only Retrieval}

The limitations of keyword-based search in cultural heritage contexts are well documented. Standard
information retrieval models assume that a user can formulate a query that matches terms in the
collection's metadata. But cultural and artistic collections often do not conform to this: their
meaning is layered, historically contingent, visually encoded, and medically or symbolically
specific in ways that metadata fields rarely capture \citep{shiri2015}.

ArtSeek, a recent multimodal framework for artwork analysis, identifies this same gap as a central
design challenge: most digitised collections lack structured links to external knowledge bases like
Wikidata or Wikipedia, meaning that conventional entity-based search cannot function at all
\citep{fanelli2025}. A system that only indexes what is explicitly labelled will logically fail to
surface what is implicitly meaningful.

The Bagatelle system prior to this project already acknowledged this limitation. Its starting
architecture included LLM-generated textual descriptions of artworks and multimodal embeddings,
a deliberate step away from raw metadata retrieval and toward semantic representation. The goal
of this thesis is to build on that foundation, not just to improve retrieval, but to change the mode
of interaction from search to exploration.

\subsection{Multimodal and Semantic Representations of Artworks}

A significant body of recent work has established that artworks are best represented by a
combination of visual and textual signals rather than by a single feature space. \citet{castellano2022}
developed ArtGraph, a knowledge graph linking artworks to artists, styles, and periods. Their
research demonstrated that combining visual embeddings from a Vision Transformer with
graph-structural embeddings from the knowledge graph outperforms purely visual approaches in
style and genre classification. The key insight is that semantic context (what an artwork is about,
who made it, in what tradition) is information that visual features alone cannot encode.

ArtSeek extends this further by integrating retrieval-augmented generation: a multimodal large
language model retrieves contextual knowledge fragments about an input artwork and uses those
to reason about visual motifs, historical context, and artistic meaning \citep{fanelli2025}. This
approach is important for the Bagatelle context because it explicitly targets collections that lack
external links, exactly the situation faced here.

A particularly relevant recent finding comes from \citet{zhang2026}, who study multi-modal
knowledge graph enrichment and show that visually complex or abstract images, which they term
``rich-semantic'', are among the hardest cases for standard visual encoders: the embeddings often
fail to capture the relevant semantics, leading to significant information loss. Converting those
images to LLM-generated text descriptions substantially recovers their semantic content. This
finding directly supports the design rationale of this project: artworks in the Bagatelle collection
are precisely this kind of rich-semantic image, and using LLM-generated descriptions alongside
CLIP embeddings is not just a convenience but a principled response to the limits of visual-only
representations.

In this project, the Bagatelle server provides both image-based CLIP embeddings (in three
variants: CLIP ViT-L/14, OpenCLIP ViT-L-14, and SigLIP~2) and LLM-generated textual
descriptions from three models (GPT-4o, Claude, and GPT-5). These form the multimodal
foundation on which the graph-based exploration layer is built. The choice of which representation
to use, visual, textual, or combined, becomes a research question in itself: different representations
surface different kinds of relationships between artworks.

\subsection{Knowledge Graphs as an Exploration Structure}

A knowledge graph represents entities as nodes and relationships as typed, directed edges. This
structure does something that a flat vector database cannot: it makes the nature of connections
explicit. Rather than just knowing that two artworks are close in embedding space, a knowledge
graph can record that they share a depicted disease, a historical period, an artistic technique, or
a symbolic motif.

Prior work on knowledge graphs in cultural heritage shows both the promise and the complexity
of this approach. The Linked Stage Graph project (FIZ Karlsruhe) built a knowledge graph from
roughly 7,000 historical theatre photographs and their metadata, enabling exploration via linked
data and graph-based navigation. A key challenge they identified is data sparsity: even a large
collection may lack the structured metadata needed to construct a dense, navigable graph
\citep{tietz2020}. This is a problem directly relevant to the Bagatelle collection, where explicit
medical and artistic metadata is limited.

At a larger scale, work on the Europeana database has demonstrated that graph-based clustering
of art objects, using topic modelling to surface semantic groupings, can reveal associations that
no individual metadata field would expose \citep{kouretsis2022}. For the purposes of this project,
the graph is not constructed from pre-existing metadata alone, but generated dynamically from
embedding similarity and LLM-derived semantic relationships, making the schema itself a design
decision.

\citet{huang2023} address the challenge of fragmented cultural heritage data directly: they propose
a joint entity-relation extraction model to build knowledge graphs from sparse, heterogeneous
museum records, showing that knowledge completion techniques can improve both the density
and usefulness of such graphs. This provides methodological context for the design choices made
in this thesis, even though the approach taken here uses embedding-based rather than
ontology-based relationship construction.

\subsection{Interpretability and Exploratory Search}

A recurring theme across this literature is that retrieval and exploration are different goals, and
that systems designed only for retrieval often fail at exploration. Retrieval optimises for returning
the most relevant result to a given query. Exploration optimises for discovering connections that
the user did not know to look for.

The Linked Stage Graph work explicitly argues that for a broad range of users to engage
meaningfully with cultural heritage data, multiple dimensions of the collection must be navigable
simultaneously \citep{tietz2020}. ArtSeek frames this as an interpretability challenge: a system
should not only retrieve related artworks but enable the user to understand \emph{why} they are
related \citep{fanelli2025}. This distinction between showing a result and explaining a connection
shapes the design of the system built in this thesis.

In the current implementation, each graph edge carries both a numerical similarity score and an
LLM-generated explanation of the semantic relationship between two artworks. This makes the
graph inspectable: a user navigating from a surgical scene in the nineteenth century to a plague
allegory from the Renaissance can read why the system considers them connected, rather than
simply trusting that it does. This aligns with the growing interest in explainable AI for cultural
heritage applications, where transparency in retrieval rationale directly affects user trust and
educational value.

\subsection{Comparing Relation Strategies}

A central contribution of this project is not just building a graph-based exploration system, but
building one that supports multiple, configurable strategies for generating graph connections.
Prior work has mostly treated the relationship between artworks as a fixed property, determined
by metadata, by a single embedding model, or by a pre-defined ontology. This project takes the
position that the right strategy depends on the exploration goal.

The ArtSeek framework provides the closest case in point for this kind of multimodal,
retrieval-oriented flexibility \citep{fanelli2025}, though it operates at the level of individual
queries rather than graph construction. \citet{zhang2026} show that text descriptions derived from
images can outperform raw visual embeddings for semantically complex artworks, a finding that
informs the design of the text-based and combined strategies in this project.

In this project, five graph generation strategies have been designed, implemented, and compared:
text-based (using LLM-generated descriptions), image-based (using CLIP visual embeddings),
combined (fusing both), query-targeted (anchoring expansion to the user's original search intent),
and LLM-confirmed (using the language model to judge the validity of candidate edges). The
evaluation design draws on the observation that user preference and expert judgment are the most
meaningful measures of exploration quality in a collection where there is no single correct answer.

\subsection{Thesis Structure}

This thesis proceeds as follows. Chapter~2 provides background on the relevant literature and
positions this project within it. Chapter~3 describes the design of the extended Bagatelle system,
including the graph schema, relation strategy architecture, and interface design. Chapter~4 covers
the technical implementation. Chapter~5 presents the evaluation design and, once the study is
complete, results. Chapter~6 discusses findings, limitations, and future directions. Chapter~7
concludes.

% ══════════════════════════════════════════════════════════════════════════════
\section{Background and Related Work}

This chapter situates the project within relevant prior work across four areas: cultural heritage
exploration systems, semantic and multimodal representations of artworks, knowledge graph
construction in the cultural domain, and evaluation of exploratory search.

\subsection{Cultural Heritage Digital Libraries and Exploration Systems}

[Overview of the shift from static digital libraries to interactive exploration systems. Discussion
of key systems including Europeana, and challenges of scale, metadata sparsity, and user
diversity.]

\subsection{Keyword-Based Retrieval and Its Limits in Art Collections}

[Detailed treatment of why standard IR models fail for visually and symbolically dense art
collections. The missing link between indexable metadata and embodied visual meaning.]

\subsection{Multimodal Representations: CLIP, LLMs, and Art}

[Technical background on CLIP (image-text alignment), its strengths for cross-modal retrieval and
weaknesses for pure text similarity. LLM-generated descriptions as an alternative textual signal.
MiniLM as a text similarity model. Discussion of the Bagatelle server's existing embedding
pipeline as the starting point. Recent work by \citet{zhang2026} on the limits of visual encoders
for rich-semantic images and the recovery of semantic content through text generation.]

\subsection{Knowledge Graphs in Cultural Heritage}

[Overview of knowledge graph construction approaches: metadata-driven, ontology-based,
clustering-based, and deep-learning-assisted.]

\subsection{Graph-Based Exploration Interfaces}

[How graph structures can be exposed to users. Layered layouts, node and edge labelling,
interactive expansion. Tension between exploration breadth and comprehensibility.]

\subsection{Evaluation of Exploratory and Semantic Search}

[Challenges of evaluating systems where there is no single correct answer. Side-by-side preference
comparison, edge relevance rating, and expert judgment as evaluation approaches.]

% ══════════════════════════════════════════════════════════════════════════════
\section{System Design}

This chapter describes the overall architecture of the extended Bagatelle system and the design
decisions behind it.

\subsection{Starting Point: The Bagatelle Server}

[Brief description of the existing system: Flask backend, Qdrant vector database, CLIP and MiniLM
embeddings, LLM-generated HTML descriptions (GPT-4o, Claude, GPT-5), gallery interface. What
was working and what was not before the project began.]

\subsection{Design Goals and Constraints}

The core design goal is to turn a search-results interface into a graph-based exploration tool. Key
constraints include staying compatible with the existing Qdrant infrastructure, avoiding the
introduction of a separate graph database at this stage, and keeping the prototype usable. The
design decision was made to build graph state in the JavaScript frontend and use Qdrant as the
semantic neighbour engine.

\subsection{Graph Schema}

Artworks are represented as nodes and relationships as edges. Information stored per node
includes: image path, title, category, LLM description, and embedding. Information stored per
edge includes: similarity score, source modality, and LLM-generated explanation. A static graph
database (e.g.\ Neo4j) was not used at this stage.

\subsection{Graph Generation Strategies}

Five strategies have been implemented, each targeting a different kind of relationship between
artworks.

\subsubsection{Text-Based Graph}
Semantic similarity between LLM-generated descriptions using MiniLM embeddings. Three
description sources are available: GPT-4o, Claude, and GPT-5.

\subsubsection{Image-Based Graph}
Visual similarity via image embeddings. Three model variants are supported: CLIP ViT-L/14,
OpenCLIP ViT-L-14, and SigLIP~2 so400m.

\subsubsection{Combined Graph}
Fused scoring from text and image retrieval. The relative weight of each modality is controlled
through a slider in the interface.

\subsubsection{Query-Targeted Graph}
Graph expansion anchored to the user's original search query, not just the clicked artwork's
description. This keeps exploration relevant to the user's starting intent across multiple hops.

\subsubsection{LLM-Confirmed Graph}
Candidate edges are generated by embedding similarity, then validated by a language model.
Each edge receives a relevance judgment (yes / maybe / no), a relation type, and a natural-language
explanation. Edges judged ``no'' are shown in the results panel but not added to the graph.

\subsection{Exclusion and Novelty in Graph Expansion}

The \texttt{exclude\_paths} mechanism prevents the graph from looping over already-seen
artworks. The rationale is that search relevance and exploration novelty are distinct goals, and
a pure nearest-neighbour strategy causes stalemates in a small collection.

\subsection{Interface Design}

The interface consists of a gallery view with clickable artworks, a graph view (layered
left-to-right layout, visible edge scores, hoverable explanations), related artwork cards with
expandable descriptions, and a clear/reset graph function. The design principle is that every
connection should be inspectable, not just visible.

% ══════════════════════════════════════════════════════════════════════════════
\section{Implementation}

This chapter documents the technical implementation of the system, covering backend
development, vector database management, and frontend graph rendering.

\subsection{Backend Architecture}

[Flask application structure. Key routes: \texttt{/retrieve} (initial search), \texttt{/related}
(graph expansion). Request/response formats.]

\subsection{Vector Database: Qdrant Setup and Schema}

[Collection structure, named vectors, payload indexing. Issues encountered: model mismatch
between CLIP and MiniLM, schema errors, client API version changes
(\texttt{client.search} $\rightarrow$ \texttt{client.query\_points}). Data pipeline: HTML cleaning
with BeautifulSoup, embedding generation, payload construction with title priority logic
(CSV metadata $\rightarrow$ HTML title $\rightarrow$ LLM-derived $\rightarrow$ filename
fallback).]

\subsection{Embedding Pipeline}

[Build scripts for text and image collections. Image path resolution and extension handling
(\texttt{.jpg}, \texttt{.jpeg}, \texttt{.png}, \texttt{.webp}, \texttt{.jfif}, \texttt{.gif},
\texttt{.avif}). Missing image handling: non-fatal warnings, skip with logging.]

\subsection{LLM Integration}

[Claude and OpenAI API integration for edge explanation generation. Prompt design for relation
explanation. Extension to LLM-confirmed mode: the model returns a relevance judgment, relation
type, and explanation.]

\subsection{Frontend Graph Engine}

[JavaScript graph state: \texttt{graphNodes} (Map), \texttt{graphEdges} (Map),
\texttt{graphRoot}. Layered depth-based layout algorithm. SVG edge rendering with invisible
hitboxes for hover. Node labels from title payload. Edge score labels. Exclusion path tracking
and transmission to backend.]

\subsection{Development Decisions and Lessons Learned}

[Embedding model choice matters just as much as architecture: CLIP is not suitable for
text-to-text similarity. Payload consistency (image paths must match static files). API sensitivity
to \texttt{qdrant-client} version. Separating category metadata from artwork titles. Full-text
storage for \texttt{section\_text} improves semantic precision.]

% ══════════════════════════════════════════════════════════════════════════════
\section{Evaluation}

This chapter describes the evaluation design for the five graph generation strategies and will
present results once the study is complete.

\subsection{Evaluation Design}

The core question the evaluation tries to answer is: which graph generation strategy produces
the most meaningful and navigable exploration paths through the Bagatelle collection? Because
there is no single objectively correct answer for exploratory search, unlike document retrieval
where relevance can be ground-truthed, the evaluation relies on human judgment rather than
automated metrics.

Three complementary evaluation components are planned, each targeting a different level of
the system:

\begin{itemize}
    \item \textbf{Side-by-side strategy comparison:} two graph modes are shown simultaneously,
    starting from the same artwork. Participants rate which side produces more relevant
    connections, which they would continue exploring, and which side's explanations are more
    illuminating.
    \item \textbf{Edge relevance judgments:} for a sample of edges from each strategy, expert
    raters assess whether the connection is art-historically, medically, or educationally
    meaningful.
    \item \textbf{Graph-level statistics:} quantitative comparison across strategies, including
    average similarity score, score variance, category distribution of connected nodes, and
    expansion depth before the graph stalls.
\end{itemize}

An MD with direct connections to the medical humanities community has agreed to participate as
an expert evaluator. This provides domain-grounded feedback on whether the graph strategies
surface connections that are meaningful from a medical and art-historical perspective, not just
mathematically close in embedding space. The evaluation interface to support this study is
currently being implemented.

\subsection{Side-by-Side Strategy Comparison}

[Two graph modes presented simultaneously from the same starting artwork. Users or expert
evaluators rate: which side produces more relevant connections; which side they would continue
exploring; which side's explanations are more illuminating. Rating interface design and question
wording.]

\subsection{Edge Relevance Judgments}

[For a sample of edges from each strategy: expert rating of whether the connection is
art-historically, medically, or educationally meaningful. Scale and rating criteria.]

\subsection{Graph-Level Statistics}

Quantitative metrics will be collected across all five strategies to complement the subjective
evaluation. These include average and minimum/maximum similarity scores per strategy, score
variance (a proxy for how selective the strategy is), category distribution of connected nodes
(whether the graph stays within one medical theme or crosses categories), and expansion depth
before the graph stalls due to the exclusion mechanism.

The working hypothesis is that the text-based graph will show higher thematic coherence within
medical categories, the image-based graph will produce more visually diverse but categorically
varied connections, the combined graph will strike a balance between the two, and the
LLM-confirmed graph will have the highest precision at the cost of coverage.

\subsection{Results}

[To be completed once the evaluation study is conducted.]

% ══════════════════════════════════════════════════════════════════════════════
\section{Discussion}

This chapter reflects on what the results reveal about graph-based exploration of art collections,
and what remains open. Sections~6.1--6.3 will be expanded once the evaluation data is available.

\subsection{What the Comparison Reveals}

[Discussion of the differences between text, image, and combined strategies in practice. Why
visual and semantic connections diverge for medical art. The role of LLM-generated descriptions
in shaping the graph structure.]

\subsection{The Role of LLM-Generated Explanations}

[Do edge explanations improve user trust and navigation quality? Are there cases where the
explanation is misleading or generic? Observations from development and testing.]

\subsection{Search Relevance vs.\ Exploration Novelty}

[The tension surfaced during development: the most similar artwork is not always the most
valuable next step in exploration. Discussion of the exclusion mechanism and whether diversity
ranking would be a better long-term approach.]

\subsection{Limitations}

Several limitations are worth noting. The Bagatelle collection of roughly 600 artworks is
relatively small, and findings about which strategy works best may not generalise to larger or
differently structured collections. The preference-based evaluation is inherently subjective;
different evaluators with different disciplinary backgrounds may prefer different strategies.
The graph currently lives only in the browser session and is not persisted, which limits how
long or complex an exploration path can be built in practice. Finally, LLM explanations are
generated per-pair and may not be fully consistent across sessions, since the models are not
deterministic.

\subsection{Future Directions}

Several extensions are worth pursuing beyond this project. The most immediately useful would
be enriching the graph with external knowledge sources: linking artworks to Wikidata entries
for depicted conditions, medical concepts from standard ontologies, or biographical data about
artists. This would allow the graph to surface connections that the collection's own descriptions
cannot, since they are limited to what the LLM infers from the image itself.

Persistent graph storage would allow users to save and return to exploration sessions, and to
share interesting paths with others. Currently the graph resets when the browser is closed.
A saved-session model would also make longitudinal evaluation possible.

On the evaluation side, a broader user study beyond expert evaluation would add useful signal:
students, general museum visitors, and medical professionals likely prioritise different kinds of
connections, and understanding those preferences would help in deciding which strategy to
recommend for which context. Finally, diversity-aware ranking, where the graph actively seeks
out artworks from underrepresented categories rather than always picking the nearest neighbour,
could address the stalemate problem more elegantly than the current exclusion mechanism.

% ══════════════════════════════════════════════════════════════════════════════
\section{Conclusion}

This project set out to turn a search-results interface into a semantically navigable graph, and
the core of that goal has been achieved. The extended Bagatelle system supports five distinct
graph generation strategies, each of which produces a different view of the relationships within
the collection. Users can start from any artwork in the catalogue, build a graph by clicking
through connections, inspect the explanation behind every edge, and switch between text-based,
image-based, combined, query-targeted, and LLM-confirmed modes to explore how different
representations shape what gets discovered.

The key design principle throughout has been that connections should be inspectable, not just
visible. A system that shows two artworks as related without explaining why places all
interpretive weight on the user. By attaching LLM-generated explanations to every edge, the
graph gives users a reason to follow or skip a connection, which is closer to how meaningful
exploration actually works.

The broader implication is that for collections where meaning is visually encoded and medically
or historically specific, as it is in the Bagatelle collection, graph-based exploration with
interpretable edges is a more appropriate interaction model than keyword search. Whether the
text strategy, the image strategy, or the combined approach best serves domain experts in
practice is the question the upcoming evaluation is designed to answer.

% ══════════════════════════════════════════════════════════════════════════════
\newpage
\bibliographystyle{apalike}
\begin{thebibliography}{99}

\bibitem[Ardissono et al.(2016)]{ardissono2016}
Ardissono, L., Lucenteforte, M., Mauro, N., Savoca, A., Voghera, A., \& La Riccia, L. (2016).
Exploration of cultural heritage information via textual search queries. In
\textit{Proceedings of the 18th International Conference on Human-Computer Interaction with
Mobile Devices and Services Adjunct} (pp. 992--1001).

\bibitem[Castellano et al.(2022)]{castellano2022}
Castellano, G., Digeno, V., Sansaro, G., \& Vessio, G. (2022). Leveraging knowledge graphs and
deep learning for automatic art analysis. \textit{Knowledge-Based Systems}, \textit{248}, 108859.

\bibitem[Fanelli et al.(2025)]{fanelli2025}
Fanelli, N., et al. (2025). ArtSeek: Deep artwork understanding via multimodal in-context
reasoning and late interaction retrieval. \textit{arXiv:2507.21917}.

\bibitem[Huang et al.(2023)]{huang2023}
Huang, Y., Yu, S., Chu, J., Fan, H., \& Du, B. (2023). Using knowledge graphs and deep learning
algorithms to enhance digital cultural heritage management. \textit{Heritage Science},
\textit{11}(1), 204.

\bibitem[Kouretsis et al.(2022)]{kouretsis2022}
Kouretsis, A., Varlamis, I., Limniati, L., Pergantis, M., \& Giannakoulopoulos, A. (2022).
Mapping art to a knowledge graph: Using data for exploring the relations among visual objects
in Renaissance art. \textit{Future Internet}, \textit{14}(7), 206.

\bibitem[Shiri et al.(2015)]{shiri2015}
Shiri, A., Ruthven, I., \& Chowdhury, G.~G. (2015). Semantic access and exploration in cultural
heritage digital libraries. In \textit{Cultural Heritage Information: Access and Management}
(pp. 177--196). Facet.

\bibitem[Tietz et al.(2020)]{tietz2020}
Tietz, T., Waitelonis, J., Alam, M., \& Sack, H. (2020). Knowledge graph based analysis and
exploration of historical theatre photographs. In \textit{Qurator}.

\bibitem[Xu(2020)]{xu2020}
Xu, M. (2020). Understanding graph embedding methods and their applications.
\textit{arXiv}. \url{https://arxiv.org/abs/2012.08019}

\bibitem[Zhang et al.(2026)]{zhang2026}
Zhang, P., Zaporojets, K., Liu, J., Huang, J.-H., \& Groth, P. (2026). Are a thousand words
better than a single picture? Beyond images: a framework for multi-modal knowledge graph
dataset enrichment. \textit{arXiv:2603.16974}.

\end{thebibliography}
\end{document}
